% ============================================================================
%  PRECÁLCULO — Sesión 14
%  Semana 14: 9 al 14 de noviembre de 2026
%  3.3 Factorización
%  Duración prevista: 3 horas
% ============================================================================
\documentclass[aspectratio=169,10pt]{beamer}

\usepackage[utf8]{inputenc}
\IfFileExists{spanish.ldf}{\usepackage[spanish,es-noshorthands]{babel}}{}
\usepackage{amsmath,amssymb}
\usepackage{booktabs}
\usepackage{tikz}

\usetheme{Uniminuto}

\uniasignatura{Precálculo}
\uniprofesor{Andrés Felipe Puerta Velez}
\unicorreo{andres.puerta-v@uniminuto.edu.co}
\uninrc{MATE1161}
\unisemana{Sesión 14 — Semana 14 · 9 al 14 de noviembre de 2026}

\newcommand{\listatiny}{\setbeamerfont{itemize/enumerate body}{size=\tiny}\tiny}

\begin{document}

\portada

\begin{frame}{Orden del día}
  \begin{enumerate}
    \item Llamado a lista y revisión de los ejercicios propuestos
    \item Repaso: los productos notables, leídos al revés
    \item \azul{Tema 3.3: los siete casos de factorización}
    \item La estrategia: en qué orden intentar
    \item Factorización completa y resolución de ecuaciones
    \item Taller, errores frecuentes y cierre
  \end{enumerate}

  \vspace{0.25cm}
  {\footnotesize\color{uniGris} Sesión de 3 horas. Factorizar es
  \destacar{la habilidad algebraica más rentable del curso}: sin ella no se
  simplifican fracciones, no se resuelven ecuaciones y no se calculan límites.}
\end{frame}

\begin{frame}{Repaso: factorizar es desandar el camino}
  \scriptsize
  \begin{center}
  \begin{tabular}{@{}p{5.4cm}p{1.4cm}p{5.0cm}@{}}
    \toprule
    \textbf{Desarrollar (sesión 13)} & & \textbf{Factorizar (hoy)} \\
    \midrule
    $(x+5)(x-5) = x^{2}-25$ & $\longleftrightarrow$ &
      $x^{2}-25 = (x+5)(x-5)$ \\
    $(x+3)^{2} = x^{2}+6x+9$ & $\longleftrightarrow$ &
      $x^{2}+6x+9 = (x+3)^{2}$ \\
    $2(x+4) = 2x+8$ & $\longleftrightarrow$ &
      $2x+8 = 2(x+4)$ \\
    \bottomrule
  \end{tabular}
  \end{center}

  \vspace{0.12cm}
  \begin{alertblock}{Qué significa factorizar}
    \scriptsize
    \destacar{Escribir una suma como un producto.} Y sirve por una razón
    concreta: \azul{con un producto se puede simplificar y se puede igualar a
    cero}; con una suma, no.
  \end{alertblock}

  \vspace{0.1cm}
  \begin{exampleblock}{La propiedad que lo justifica todo}
    \scriptsize
    \[ a\cdot b = 0 \quad\Longleftrightarrow\quad a = 0 \ \text{ o } \ b = 0 \]
    {\tiny\color{uniGris} No existe nada parecido para la suma: de $a+b=0$ no
    se deduce nada sobre $a$ ni sobre $b$ por separado. \destacar{Por eso hay
    que factorizar antes de resolver.}}
  \end{exampleblock}
\end{frame}

\seccion{3.3 Factorización}

\begin{frame}{Casos 1 y 2 — factor común}
  \scriptsize
  \begin{columns}[T,onlytextwidth]
    \begin{column}{0.48\textwidth}
      \begin{block}{Factor común monomio}
        \tiny
        Se saca el \azul{MCD de los coeficientes} y la letra con el
        \azul{menor exponente}:
        \[ 6x^{3}-9x^{2} = \destacar{3x^{2}(2x-3)} \]
        \[ 4a^{2}b+8ab^{2} = \destacar{4ab(a+2b)} \]
      \end{block}
    \end{column}
    \begin{column}{0.48\textwidth}
      \begin{block}{Factor común por agrupación}
        \tiny
        Cuatro términos: se agrupan de a dos.
        \[ ax+ay+bx+by \]
        \[ = a(x+y) + b(x+y) \]
        \[ = \destacar{(a+b)(x+y)} \]
      \end{block}
    \end{column}
  \end{columns}

  \vspace{0.12cm}
  \begin{alertblock}{La señal de que la agrupación va bien}
    \scriptsize
    Después de sacar factor común en cada pareja, \destacar{los dos paréntesis
    deben quedar idénticos}. Si no quedan iguales, hay que
    \azul{reagrupar de otra forma} o el polinomio no se factoriza por este
    caso.
  \end{alertblock}

  \vspace{0.05cm}
  {\tiny\color{uniGris} \destacar{El factor común siempre se intenta primero.}
  Sacarlo suele dejar adentro un polinomio más simple, que a veces se factoriza
  otra vez.}
\end{frame}

\begin{frame}{Casos 3 y 4 — los que vienen de productos notables}
  \scriptsize
  \begin{columns}[T,onlytextwidth]
    \begin{column}{0.48\textwidth}
      \begin{block}{Diferencia de cuadrados}
        \tiny
        \[ a^{2}-b^{2} = (a+b)(a-b) \]
        \[ x^{2}-49 = \destacar{(x+7)(x-7)} \]
        \[ 4a^{2}-25b^{2} = \destacar{(2a+5b)(2a-5b)} \]
        \azul{Solo con resta.} $x^2+49$ no se factoriza en $\mathbb{R}$.
      \end{block}
    \end{column}
    \begin{column}{0.48\textwidth}
      \begin{block}{Trinomio cuadrado perfecto}
        \tiny
        \[ a^{2}\pm 2ab+b^{2} = (a\pm b)^{2} \]
        \[ x^{2}+10x+25 = \destacar{(x+5)^{2}} \]
        \[ 9y^{2}-24y+16 = \destacar{(3y-4)^{2}} \]
      \end{block}
    \end{column}
  \end{columns}

  \vspace{0.1cm}
  \begin{alertblock}{Cómo reconocer un trinomio cuadrado perfecto}
    \scriptsize
    \textbf{1.} El primero y el último son \azul{cuadrados perfectos}. \quad
    \textbf{2.} El del medio es \destacar{el doble producto} de sus raíces.
    \vspace{0.06cm}
    En $9y^{2}-24y+16$: raíces $3y$ y $4$; doble producto
    $2(3y)(4) = 24y$ \checkmark. \azul{Si no cumple el punto 2, no es
    perfecto} y hay que ir al caso 5.
  \end{alertblock}
\end{frame}

\begin{frame}{Casos 5 y 6 — trinomios de segundo grado}
  \scriptsize
  \begin{columns}[T,onlytextwidth]
    \begin{column}{0.46\textwidth}
      \begin{block}{$x^{2}+bx+c$}
        \tiny
        Se buscan dos números que \azul{multiplicados den $c$} y
        \azul{sumados den $b$}.
        \[ x^{2}+7x+12 = \destacar{(x+3)(x+4)} \]
        \[ x^{2}-x-20 = \destacar{(x-5)(x+4)} \]
        {\color{uniGris} $-5\cdot 4 = -20$ y $-5+4 = -1$ \checkmark}
      \end{block}
    \end{column}
    \begin{column}{0.50\textwidth}
      \begin{block}{$ax^{2}+bx+c$ — método del producto}
        \tiny
        En $2x^{2}+7x+3$: se buscan dos números que den
        $a\cdot c = 6$ y sumen $7$: son $6$ y $1$.
        \[ 2x^{2}+6x+x+3 \]
        \[ = 2x(x+3) + 1(x+3) = \destacar{(2x+1)(x+3)} \]
      \end{block}
    \end{column}
  \end{columns}

  \vspace{0.1cm}
  \begin{exampleblock}{Otro con $a\neq 1$}
    \scriptsize
    $6x^{2}-x-2$: \ $a\cdot c = -12$; dos números que den $-12$ y sumen $-1$:
    $-4$ y $3$.
    \[ 6x^{2}-4x+3x-2 = 2x(3x-2)+1(3x-2) = \destacar{(2x+1)(3x-2)} \]
  \end{exampleblock}

  \vspace{0.05cm}
  {\tiny\color{uniGris} El método es siempre el mismo: \destacar{partir el
  término del medio} usando el producto $a\cdot c$, y después agrupar. Reduce
  el caso 6 al caso 2.}
\end{frame}

\begin{frame}{Caso 7 — suma y diferencia de cubos}
  \scriptsize
  \begin{columns}[T,onlytextwidth]
    \begin{column}{0.48\textwidth}
      \begin{block}{Las dos fórmulas}
        \tiny
        \[ a^{3}+b^{3} = (a+b)(a^{2}-ab+b^{2}) \]
        \[ a^{3}-b^{3} = (a-b)(a^{2}+ab+b^{2}) \]
        \[ x^{3}+27 = \destacar{(x+3)(x^{2}-3x+9)} \]
        \[ 8x^{3}-1 = \destacar{(2x-1)(4x^{2}+2x+1)} \]
      \end{block}
    \end{column}
    \begin{column}{0.48\textwidth}
      \begin{alertblock}{Cómo recordar los signos}
        \tiny
        \destacar{El binomio lleva el mismo signo} que la expresión original.
        \destacar{El trinomio lleva el signo contrario en el término del
        medio} y \azul{siempre} $+$ al final.

        \vspace{0.08cm}
        Y el trinomio \azul{nunca} es cuadrado perfecto: le falta el doble
        producto (tiene $ab$, no $2ab$).
      \end{alertblock}
    \end{column}
  \end{columns}

  \vspace{0.1cm}
  \begin{exampleblock}{A diferencia de los cuadrados\dots}
    \scriptsize
    La \destacar{suma} de cubos \azul{sí} se factoriza, mientras que la suma de
    cuadrados no. Compare: $x^{2}+9$ es irreducible en $\mathbb{R}$, pero
    $x^{3}+27$ no lo es.
  \end{exampleblock}
\end{frame}

\begin{frame}{La estrategia: en qué orden intentar}
  \scriptsize
  \begin{center}
  \begin{tabular}{@{}cp{4.4cm}p{5.8cm}@{}}
    \toprule
    & \textbf{Se pregunta} & \textbf{Si la respuesta es sí} \\
    \midrule
    1 & ¿Hay factor común? &
      Sáquelo \azul{siempre primero} y siga con lo de adentro \\
    2 & ¿Cuántos términos quedan? & \\
    3 & \quad Dos: ¿resta de cuadrados? ¿suma o resta de cubos? &
      Casos 3 y 7 \\
    4 & \quad Tres: ¿cuadrado perfecto? Si no, trinomio &
      Casos 4, 5 y 6 \\
    5 & \quad Cuatro: agrupar de a dos & Caso 2 \\
    6 & ¿Algún factor se puede volver a factorizar? &
      \destacar{Repita hasta que no se pueda más} \\
    \bottomrule
  \end{tabular}
  \end{center}

  \vspace{0.08cm}
  \begin{exampleblock}{Factorización completa: un ejemplo con tres pasos}
    \tiny
    \[ 3x^{3}-27x
       \;\xrightarrow{\text{factor común}}\; 3x\left(x^{2}-9\right)
       \;\xrightarrow{\text{dif.\ de cuadrados}}\;
       \destacar{3x(x+3)(x-3)} \]
    Quedarse en $3x(x^2-9)$ es \azul{un ejercicio a medias}.
  \end{exampleblock}
\end{frame}

\begin{frame}{Ejercicio 1 — en clase}
  \begin{alertblock}{Individual (15 minutos)}
    \scriptsize
    Factorice, indicando qué caso usó:

    \vspace{0.25cm}
    \centering
    \begin{tabular}{@{}lll@{}}
      1. \ $10x^{4}-15x^{2}$ &
      2. \ $x^{2}-64$ &
      3. \ $x^{2}+12x+36$ \\[0.9em]
      4. \ $x^{2}+9x+20$ &
      5. \ $x^{2}-3x-28$ &
      6. \ $9a^{2}-16b^{2}$ \\[0.9em]
      7. \ $2x^{2}+7x+3$ &
      8. \ $x^{3}-8$ &
      9. \ $mx+my+nx+ny$ \\
    \end{tabular}
  \end{alertblock}
\end{frame}

\begin{frame}{Ejercicio 1 — solución}
  \scriptsize
  \begin{center}
  \begin{tabular}{@{}clp{3.0cm}@{\hspace{0.5cm}}clp{2.8cm}@{}}
    \toprule
    \textbf{\#} & \textbf{Factorización} & \textbf{Caso} &
    \textbf{\#} & \textbf{Factorización} & \textbf{Caso} \\
    \midrule
    1 & $5x^{2}(2x^{2}-3)$ & Factor común &
    6 & $(3a+4b)(3a-4b)$ & Dif.\ cuadrados \\[0.25em]
    2 & $(x+8)(x-8)$ & Dif.\ cuadrados &
    7 & $(2x+1)(x+3)$ & Trinomio $a\neq1$ \\[0.25em]
    3 & $(x+6)^{2}$ & Cuadrado perfecto &
    8 & $(x-2)(x^{2}+2x+4)$ & Dif.\ de cubos \\[0.25em]
    4 & $(x+4)(x+5)$ & Trinomio simple &
    9 & $(m+n)(x+y)$ & Agrupación \\[0.25em]
    5 & $(x-7)(x+4)$ & Trinomio simple & & & \\
    \bottomrule
  \end{tabular}
  \end{center}

  \vspace{0.1cm}
  \begin{alertblock}{Verificación obligatoria}
    \tiny
    \destacar{Multiplique de vuelta.} En el 5: $(x-7)(x+4) = x^2+4x-7x-28
    = x^2-3x-28$ \checkmark. Es la única forma de estar seguro, y toma diez
    segundos.
  \end{alertblock}
\end{frame}

\begin{frame}{Ejercicio 2 — en clase}
  \begin{alertblock}{En parejas (18 minutos)}
    \scriptsize
    Factorice \destacar{completamente} (puede hacer falta más de un caso):

    \vspace{0.25cm}
    \begin{columns}[T,onlytextwidth]
      \begin{column}{0.48\textwidth}
        \begin{enumerate}\setlength{\itemsep}{0.2cm}
          \item $x^{3}-4x$
          \item $3x^{2}-27$
          \item $2x^{3}+16$
        \end{enumerate}
      \end{column}
      \begin{column}{0.48\textwidth}
        \begin{enumerate}\setlength{\itemsep}{0.2cm}
          \setcounter{enumi}{3}
          \item $x^{4}-16$
          \item $5x^{2}+30x+45$
          \item $6x^{2}-x-2$
        \end{enumerate}
      \end{column}
    \end{columns}
  \end{alertblock}
\end{frame}

\begin{frame}{Ejercicio 2 — solución}
  \tiny
  \begin{columns}[T,onlytextwidth]
    \begin{column}{0.49\textwidth}
      \begin{block}{Puntos 1 a 3}
        \tiny
        \begin{align*}
          1)\;& x(x^{2}-4) = \destacar{x(x+2)(x-2)} \\[0.3em]
          2)\;& 3(x^{2}-9) = \destacar{3(x+3)(x-3)} \\[0.3em]
          3)\;& 2(x^{3}+8) = \destacar{2(x+2)(x^{2}-2x+4)}
        \end{align*}
      \end{block}
    \end{column}
    \begin{column}{0.49\textwidth}
      \begin{block}{Puntos 4 a 6}
        \tiny
        \begin{align*}
          4)\;& (x^{2}+4)(x^{2}-4)
                = \destacar{(x^{2}+4)(x+2)(x-2)} \\[0.3em]
          5)\;& 5(x^{2}+6x+9) = \destacar{5(x+3)^{2}} \\[0.3em]
          6)\;& \destacar{(2x+1)(3x-2)}
        \end{align*}
      \end{block}
    \end{column}
  \end{columns}

  \vspace{0.08cm}
  \begin{alertblock}{Los dos detalles finos}
    \tiny
    En el \destacar{4}, el factor $x^{2}+4$ \azul{no se factoriza más} en los
    reales: es una \emph{suma} de cuadrados. \quad
    En el \destacar{1}, el $x$ que se sacó como factor común \azul{es parte de
    la respuesta}: olvidarlo cambia el resultado.
  \end{alertblock}
\end{frame}

\begin{frame}{Para qué sirve: resolver ecuaciones}
  \scriptsize
  \begin{columns}[T,onlytextwidth]
    \begin{column}{0.48\textwidth}
      \begin{exampleblock}{Procedimiento}
        \tiny
        \textbf{1.} Igualar a cero. \quad
        \textbf{2.} Factorizar. \quad
        \textbf{3.} Igualar cada factor a cero.
        \[ x^{2}-x-20 = 0 \]
        \[ (x-5)(x+4) = 0 \]
        \[ x = \destacar{5} \quad\text{o}\quad x = \destacar{-4} \]
      \end{exampleblock}
    \end{column}
    \begin{column}{0.48\textwidth}
      \begin{exampleblock}{Otro caso}
        \tiny
        \[ x^{3} = 4x \;\Rightarrow\; x^{3}-4x = 0 \]
        \[ x(x+2)(x-2) = 0 \]
        \[ x = \destacar{0},\ \destacar{-2},\ \destacar{2} \]
        {\color{uniGris} Tres soluciones, porque el grado es 3.}
      \end{exampleblock}
    \end{column}
  \end{columns}

  \vspace{0.08cm}
  \begin{alertblock}{El error que cuesta una solución}
    \scriptsize
    En $x^{3}=4x$, \destacar{no se puede dividir entre $x$}: eso eliminaría la
    solución $x=0$. \azul{Nunca se divide entre una expresión que puede valer
    cero:} se pasa todo a un lado y se factoriza.
  \end{alertblock}
\end{frame}

\begin{frame}{Ejercicio 3 — en clase (taller principal)}
  \begin{alertblock}{En grupos de 3 (20 minutos)}
    \tiny
    \setbeamerfont{itemize/enumerate body}{size=\tiny}
    \begin{enumerate}\setlength{\itemsep}{0.05cm}
      \item Factorice completamente: \ $4x^{3}-36x$, \ $x^{2}-14x+49$, \
            $27+x^{3}$, \ $8x^{2}+2x-3$.
      \item Resuelva factorizando: \ $x^{2}+2x-15 = 0$ \ y \
            $2x^{2}-8 = 0$.
      \item Resuelva: \ $x^{3} = 9x$. ¿Cuántas soluciones tiene?
      \item El área de un rectángulo es $x^{2}+9x+20$. Escriba sus dos lados
            y calcule el perímetro para $x = 6$.
      \item Encuentre el error: \ $x^{2}-16 = (x-4)^{2}$
      \item Encuentre el error: \ ``de $x^{2}=5x$ se sigue que $x=5$''
    \end{enumerate}
  \end{alertblock}
\end{frame}

\begin{frame}{Ejercicio 3 — solución}
  \scriptsize
  \begin{columns}[T,onlytextwidth]
    \begin{column}{0.49\textwidth}
      \begin{block}{Puntos 1 a 4}
        \tiny
        \textbf{1.} $4x(x+3)(x-3)$; \ $(x-7)^{2}$; \
        $(x+3)(x^{2}-3x+9)$; \ $(4x+3)(2x-1)$

        \vspace{0.06cm}
        \textbf{2.} $(x+5)(x-3)=0 \Rightarrow \destacar{x=-5,\ 3}$; \\
        $2(x+2)(x-2)=0 \Rightarrow \destacar{x=\pm 2}$

        \vspace{0.06cm}
        \textbf{3.} $x(x+3)(x-3)=0 \Rightarrow \destacar{x=0,\ \pm3}$:
        \azul{tres} soluciones

        \vspace{0.06cm}
        \textbf{4.} Lados $(x+4)$ y $(x+5)$; con $x=6$:
        $P = 2(10+11) = \destacar{42}$
      \end{block}
    \end{column}
    \begin{column}{0.49\textwidth}
      \begin{alertblock}{Puntos 5 y 6}
        \tiny
        \textbf{5.} $x^{2}-16 = \destacar{(x+4)(x-4)}$. En cambio
        $(x-4)^2 = x^2-8x+16$: \azul{tiene término del medio}.

        \vspace{0.08cm}
        \textbf{6.} Se dividió entre $x$ y \destacar{se perdió la solución
        $x=0$}. Lo correcto: $x^2-5x=0 \Rightarrow x(x-5)=0 \Rightarrow
        x=0$ o $x=5$.
      \end{alertblock}
    \end{column}
  \end{columns}

  \vspace{0.05cm}
  {\tiny\color{uniGris} El punto \destacar{6} es el error más costoso del
  álgebra: perder soluciones al dividir. En programación equivale a un
  \azul{caso borde} sin cubrir.}
\end{frame}

\begin{frame}{Errores frecuentes en esta sesión}
  \scriptsize
  \begin{center}
  \begin{tabular}{@{}p{5.2cm}p{6.6cm}@{}}
    \toprule
    \textbf{Error} & \textbf{Corrección} \\
    \midrule
    $x^{2}-16 = (x-4)^{2}$ &
      Es $(x+4)(x-4)$: diferencia de cuadrados \\
    \addlinespace
    Factorizar $x^{2}+9$ &
      La \azul{suma} de cuadrados es irreducible en $\mathbb{R}$ \\
    \addlinespace
    Dejar $3x(x^{2}-9)$ como respuesta &
      Falta un paso: $3x(x+3)(x-3)$ \\
    \addlinespace
    Perder el factor común en la respuesta &
      El $x$ o el $3$ que se sacó \azul{es parte del resultado} \\
    \addlinespace
    Dividir entre $x$ al resolver &
      Se pierde la solución $x=0$: hay que factorizar \\
    \addlinespace
    No verificar &
      Multiplicar de vuelta detecta el error en diez segundos \\
    \bottomrule
  \end{tabular}
  \end{center}
\end{frame}

\begin{frame}{Cierre de la sesión}
  \begin{block}{Lo que debe quedar claro hoy}
    \scriptsize
    \begin{itemize}\setlength{\itemsep}{0.02cm}
      \item Factorizar es \azul{escribir una suma como producto}: los
            productos notables al revés.
      \item Los siete casos: factor común, agrupación, diferencia de cuadrados,
            cuadrado perfecto, trinomio simple, trinomio con $a\neq1$ y cubos.
      \item \destacar{El factor común siempre se saca primero}, y después se
            cuenta cuántos términos quedan.
      \item La factorización debe ser \azul{completa}.
      \item $ab=0 \Rightarrow a=0$ o $b=0$: por eso se resuelven ecuaciones
            factorizando.
      \item \destacar{Nunca se divide entre una expresión que pueda ser cero.}
    \end{itemize}
  \end{block}

  \vspace{0.02cm}
  \begin{alertblock}{Trabajo independiente}
    \tiny
    Resolver los 20 ejercicios de la lámina siguiente. En la sesión 15
    aplicamos la factorización a las \destacar{fracciones algebraicas}, vemos
    \destacar{funciones trigonométricas} y hacemos el repaso para el Examen 3.
  \end{alertblock}
\end{frame}

\begin{frame}{Ejercicios propuestos}
  \vspace{-0.1cm}
  \begin{beamercolorbox}[wd=\textwidth,sep=0.10cm]{block body}
    \centering\scriptsize
    \textbf{1--12:} factorice. \quad \textbf{13--17:} factorización completa.
    \quad \textbf{18--20:} ecuaciones.
  \end{beamercolorbox}
  \listatiny
  \vspace{0.15cm}
  \begin{columns}[T,onlytextwidth]
    \begin{column}{0.47\textwidth}
      \begin{enumerate}\setlength{\itemsep}{0.16cm}
        \item $8x^{3}+12x^{2}$
        \item $x^{2}-81$
        \item $25a^{2}-49b^{2}$
        \item $x^{2}+16x+64$
        \item $4y^{2}-20y+25$
        \item $x^{2}+11x+30$
        \item $x^{2}-2x-35$
        \item $3x^{2}+10x+8$
        \item $x^{3}+64$
        \item $27x^{3}-8$
      \end{enumerate}
    \end{column}
    \begin{column}{0.50\textwidth}
      \begin{enumerate}\setlength{\itemsep}{0.16cm}
        \setcounter{enumi}{10}
        \item $2ax-6a+bx-3b$
        \item $x^{2}-\dfrac{1}{9}$
        \item $5x^{3}-45x$
        \item $x^{4}-81$
        \item $2x^{2}+12x+18$
        \item $x^{3}-x$
        \item $4x^{3}-4x^{2}-24x$
        \item Resuelva $x^{2}-7x+12=0$
        \item Resuelva $3x^{2}-12=0$
        \item Resuelva $x^{3}=16x$
      \end{enumerate}
    \end{column}
  \end{columns}
\end{frame}

\begin{frame}{Respuestas}{Autoevaluación de los ejercicios propuestos}
  \listatiny
  \vspace{0.15cm}
  \begin{columns}[T,onlytextwidth]
    \begin{column}{0.47\textwidth}
      \begin{enumerate}\setlength{\itemsep}{0.16cm}
        \item $4x^{2}(2x+3)$
        \item $(x+9)(x-9)$
        \item $(5a+7b)(5a-7b)$
        \item $(x+8)^{2}$
        \item $(2y-5)^{2}$
        \item $(x+5)(x+6)$
        \item $(x-7)(x+5)$
        \item $(3x+4)(x+2)$
        \item $(x+4)(x^{2}-4x+16)$
        \item $(3x-2)(9x^{2}+6x+4)$
      \end{enumerate}
    \end{column}
    \begin{column}{0.50\textwidth}
      \begin{enumerate}\setlength{\itemsep}{0.16cm}
        \setcounter{enumi}{10}
        \item $(x-3)(2a+b)$
        \item $\left(x+\tfrac{1}{3}\right)\left(x-\tfrac{1}{3}\right)$
        \item $5x(x+3)(x-3)$
        \item $(x^{2}+9)(x+3)(x-3)$
        \item $2(x+3)^{2}$
        \item $x(x+1)(x-1)$
        \item $4x(x-3)(x+2)$
        \item $x=3$ o $x=4$
        \item $x=\pm 2$
        \item $x=0$, $x=4$, $x=-4$
      \end{enumerate}
    \end{column}
  \end{columns}
\end{frame}

\begin{frame}{Referencias}
  \footnotesize
  \begin{itemize}
    \item[{[1]}] Ramos del Olmo, Á. M. (2016). \textit{Matemáticas básicas
      para el acceso a la universidad}. Ediciones Pirámide.
      \texttt{https://elibro.net/es/ereader/uniminuto/49134}
    \item[{[2]}] González Pulido, J. W., Chacón Benavides, J. A. \&
      Fonseca Correa, L. Á. (2023). \textit{Matemática básica} (1.\textsuperscript{a} ed.).
      Universidad Mariana.
      \texttt{https://elibro.net/es/ereader/uniminuto/280647}
    \item[{[3]}] Placencia Valero, J. (2008). \textit{Compendio de matemática
      básica elemental} (1.\textsuperscript{a} ed.). Editorial Tébar Flores.
      \texttt{https://elibro.net/es/ereader/uniminuto/269758}
  \end{itemize}

  \vspace{0.35cm}
  \begin{center}
    {\Large\color{uniAzul}\textbf{¿Preguntas?}}
  \end{center}
\end{frame}

\end{document}
